% =====================================================================
\chapter{Soluciones de los ejercicios}
\label{ap:soluciones}
% =====================================================================

Cada solución se presenta en cinco partes: datos, planteamiento, desarrollo, resultado y comprobación. En los pasos intermedios se redondea a cuatro cifras significativas salvo que se indique otra cosa; los resultados exactos se dan como fracción o radical cuando existe.

% ---------------------------------------------------------------------
\section{Tema 1}

\begin{solucion}{eje:1-combinacion}
\datos $u=(1,2)$, $v=(3,-1)$, objetivo $(5,3)$.
\planteamiento igualar componente a componente $au+bv=(a+3b,\,2a-b)$ con $(5,3)$ y resolver el sistema.
\pasos
\begin{enumerate}[label=\arabic*.]
  \item Sistema: $a+3b=5$ y $2a-b=3$.
  \item De la segunda ecuación, $b=2a-3$. Sustituyendo en la primera: $a+3(2a-3)=5$, es decir, $7a-9=5$, luego $a=2$.
  \item $b=2\cdot 2-3=1$.
\end{enumerate}
\resultado \textbf{$a=2$, $b=1$: $(5,3)=2u+v$.} Geométricamente, $(5,3)$ es la diagonal del paralelogramo de lados $2u$ y $v$.
\comprobacion $2(1,2)+(3,-1)=(2+3,\,4-1)=(5,3)$, que coincide con el \cref{ej:operaciones}.
\end{solucion}

\begin{solucion}{eje:1-recuentos}
\datos $w\in\R^{n}$ vector de recuentos; diccionario (el, gato, come, pez); documento «el gato come el pez».
\planteamiento interpretar cada producto escalar y cada componente según la definición de $w$.
\pasos
\begin{enumerate}[label=(\alph*)]
  \item $\unos\cdot w=\sum_iw_i$ suma los recuentos de todas las palabras del diccionario: es el \textbf{número total de palabras} del documento.
  \item $w_{282}=0$ significa que la palabra 282 del diccionario \textbf{no aparece} en el documento.
  \item $h_i=w_i/(\unos\cdot w)$, es decir, $h=\dfrac{1}{\unos\cdot w}\,w$, definido si el documento tiene al menos una palabra.
  \item «el» aparece dos veces y las demás una: $w=(2,1,1,1)$; $\unos\cdot w=2+1+1+1=5$; $h=\tfrac15(2,1,1,1)=(0{,}4;\,0{,}2;\,0{,}2;\,0{,}2)$.
\end{enumerate}
\resultado \textbf{(a) número total de palabras; (b) la palabra 282 no aparece; (c) $h=w/(\unos\cdot w)$; (d) $w=(2,1,1,1)$, $\unos\cdot w=5$, $h=(0{,}4;\,0{,}2;\,0{,}2;\,0{,}2)$.}
\comprobacion $\unos\cdot h=0{,}4+0{,}2+0{,}2+0{,}2=1$: el histograma es una distribución de probabilidad sobre el diccionario.
\end{solucion}

\begin{solucion}{eje:1-normas}
\datos $x=(3,-4,12)$.
\planteamiento aplicar las definiciones de $\norm{\cdot}_1$, $\norm{\cdot}$ y $\norm{\cdot}_\infty$; después dividir por la norma euclídea.
\pasos
\begin{enumerate}[label=\arabic*.]
  \item $\norm{x}_1=\abs{3}+\abs{-4}+\abs{12}=3+4+12=19$.
  \item $\norm{x}=\sqrt{9+16+144}=\sqrt{169}=13$.
  \item $\norm{x}_\infty=\max\{3,4,12\}=12$.
  \item Orden: $12\le 13\le 19$.
  \item $\hat x=x/13=(3/13,\,-4/13,\,12/13)\approx(0{,}2308;\,-0{,}3077;\,0{,}9231)$.
\end{enumerate}
\resultado \textbf{$\norm{x}_1=19$, $\norm{x}=13$, $\norm{x}_\infty=12$; $\hat x=\frac{1}{13}(3,-4,12)$.}
\comprobacion $\norm{\hat x}^{2}=(9+16+144)/169=1$.
\end{solucion}

\begin{solucion}{eje:1-cs}
\datos $u,v\in\R^{d}$ con $v\neq 0$ (si $v=0$ ambos miembros valen~0).
\planteamiento un polinomio de segundo grado con coeficiente principal positivo que no toma valores negativos tiene discriminante $\le 0$.
\pasos
\begin{enumerate}[label=\arabic*.]
  \item $f(t)=\norm{u-tv}^{2}\ge 0$ para todo $t\in\R$, por ser una norma al cuadrado.
  \item Por la linealidad y la simetría del producto escalar, $f(t)=(u-tv)\cdot(u-tv)=\norm{v}^{2}t^{2}-2(u\cdot v)\,t+\norm{u}^{2}$.
  \item Es un polinomio $\alpha t^{2}+\beta t+\gamma$ con $\alpha=\norm{v}^{2}>0$, $\beta=-2(u\cdot v)$ y $\gamma=\norm{u}^{2}$. Si su discriminante fuera positivo, tendría dos raíces reales distintas y sería negativo entre ellas; como $f\ge 0$, debe ser $\beta^{2}-4\alpha\gamma\le 0$.
  \item Sustituyendo: $4(u\cdot v)^{2}-4\norm{u}^{2}\norm{v}^{2}\le 0$, es decir, $(u\cdot v)^{2}\le\norm{u}^{2}\norm{v}^{2}$.
  \item Tomando raíces cuadradas (ambos miembros son $\ge 0$): $\abs{u\cdot v}\le\norm{u}\norm{v}$.
  \item Si $u,v\neq 0$, dividiendo por $\norm{u}\norm{v}>0$: $\abs{u\cdot v}/(\norm{u}\norm{v})\le 1$.
\end{enumerate}
\resultado \textbf{$\abs{u\cdot v}\le\norm{u}\norm{v}$ y, por tanto, $-1\le u\cdot v/(\norm{u}\norm{v})\le 1$.}
\comprobacion con $u=(1,2)$ y $v=(3,4)$: $\abs{u\cdot v}=11$ y $\norm{u}\norm{v}=\sqrt5\cdot 5\approx 11{,}18$. La desigualdad se cumple con poco margen porque los vectores son casi paralelos.
\end{solucion}

\begin{solucion}{eje:1-paralelogramo}
\datos $u,v\in\R^{d}$; en (c), $u=(1,0)$ y $v=(0,1)$.
\planteamiento desarrollar $\norm{u\pm v}^{2}$ con la linealidad del producto escalar, y evaluar ambos miembros con la norma $\ell_1$.
\pasos
\begin{enumerate}[label=\arabic*.]
  \item $\norm{u+v}^{2}=\norm{u}^{2}+2\,u\cdot v+\norm{v}^{2}$ y $\norm{u-v}^{2}=\norm{u}^{2}-2\,u\cdot v+\norm{v}^{2}$.
  \item (a) Sumando: $\norm{u+v}^{2}+\norm{u-v}^{2}=2\norm{u}^{2}+2\norm{v}^{2}$.
  \item (b) Restando: $\norm{u+v}^{2}-\norm{u-v}^{2}=4\,u\cdot v$, de donde se obtiene la identidad de polarización.
  \item (c) $u+v=(1,1)$ y $u-v=(1,-1)$, luego $\norm{u+v}_1=\norm{u-v}_1=2$ y el primer miembro vale $2^{2}+2^{2}=8$.
  \item El segundo miembro vale $2\norm{u}_1^{2}+2\norm{v}_1^{2}=2+2=4\neq 8$.
  \item Los pasos 1 y 2 solo usan la simetría y la linealidad, válidas para todo producto interior. Si la norma $\ell_1$ fuera $\sqrt{\ip{u}{u}}$ para algún producto interior, cumpliría la ley del paralelogramo, lo que contradice el paso~5.
\end{enumerate}
\resultado \textbf{Se cumplen la ley del paralelogramo y la identidad de polarización; la norma $\ell_1$ no cumple la primera ($8\neq 4$) y, por tanto, no procede de ningún producto interior.}
\comprobacion con la norma euclídea y los mismos vectores: $\norm{u+v}^{2}+\norm{u-v}^{2}=2+2=4=2\cdot 1+2\cdot 1$; y la polarización da $\tfrac14(2-2)=0=u\cdot v$.
\end{solucion}

% ---------------------------------------------------------------------
\section{Tema 2}

\begin{solucion}{eje:2-escala}
\datos $a=(1,2)$, $b=(2,4)$, $c=(10,20)$.
\planteamiento $u\cdot v=\sum_iu_iv_i$ y $\cos(u,v)=u\cdot v/(\norm{u}\norm{v})$.
\pasos
\begin{enumerate}[label=\arabic*.]
  \item $a\cdot b=1\cdot 2+2\cdot 4=10$ y $a\cdot c=1\cdot 10+2\cdot 20=50$.
  \item $\norm{a}=\sqrt5$, $\norm{b}=\sqrt{20}$, $\norm{c}=\sqrt{500}$.
  \item $\cos(a,b)=10/\sqrt{5\cdot 20}=10/10=1$ y $\cos(a,c)=50/\sqrt{5\cdot 500}=50/50=1$.
\end{enumerate}
\resultado \textbf{$a\cdot b=10$, $a\cdot c=50$; ambos cosenos valen~1.} El producto escalar crece con la longitud de los vectores; la similitud coseno solo mide la dirección, que es la misma en los tres.
\comprobacion $c=10a$, luego $a\cdot c=10\,(a\cdot a)=10\cdot 5=50$; y, por la \cref{prop:coseno}(ii), vectores múltiplos positivos tienen coseno~1.
\end{solucion}

\begin{solucion}{eje:2-vf}
\planteamiento buscar un contraejemplo o aplicar la definición.
\pasos
\begin{enumerate}[label=(\alph*)]
  \item \textbf{Falsa.} $(1,0)\cdot(0,1)=0$ y ninguno es nulo: $u\cdot v=0$ significa que son ortogonales.
  \item \textbf{Verdadera.} $\cos(u,u)=u\cdot u/(\norm{u}\norm{u})=\norm{u}^{2}/\norm{u}^{2}=1$ si $u\neq 0$.
  \item \textbf{Falsa.} $\cos((1,2),(2,4))=1$ y los vectores son distintos. Por la \cref{prop:coseno}(ii), $\cos(u,v)=1$ solo implica $u=cv$ con $c>0$.
\end{enumerate}
\resultado \textbf{(a) falsa; (b) verdadera; (c) falsa.}
\comprobacion en (b), con $u=(3,4)$: $25/(5\cdot 5)=1$.
\end{solucion}

\begin{solucion}{eje:2-angulo}
\datos $u=(2,-1,2)$, $v=(1,2,2)$.
\planteamiento producto escalar, normas, coseno y $\theta=\arccos(\cos(u,v))$ (\cref{def:angulo}).
\pasos
\begin{enumerate}[label=\arabic*.]
  \item $u\cdot v=2\cdot 1+(-1)\cdot 2+2\cdot 2=2-2+4=4$.
  \item $\norm{u}=\sqrt{4+1+4}=3$ y $\norm{v}=\sqrt{1+4+4}=3$.
  \item $\cos(u,v)=4/(3\cdot 3)=4/9\approx 0{,}4444$.
  \item $\theta=\arccos(0{,}4444)\approx 63{,}61^\circ$.
\end{enumerate}
\resultado \textbf{$u\cdot v=4$; $\norm{u}=\norm{v}=3$; $\cos(u,v)=4/9\approx 0{,}4444$; $\theta\approx 63{,}6^\circ$.}
\comprobacion $\cos 60^\circ=0{,}5$ y $\cos 90^\circ=0$; un coseno de $0{,}44$ corresponde a un ángulo entre ambos, más próximo a $60^\circ$.
\end{solucion}

\begin{solucion}{eje:2-vecinos}
\datos $a=(1,1)$, $b=(3,3)$, $c=(1,-1)$.
\planteamiento comparar la distancia $\norm{a-x}$ con la similitud $\cos(a,x)$ para $x=b,c$.
\pasos
\begin{enumerate}[label=\arabic*.]
  \item $a-b=(-2,-2)$ y $\norm{a-b}=\sqrt8\approx 2{,}828$.
  \item $a-c=(0,2)$ y $\norm{a-c}=2$.
  \item $a\cdot b=6$, $\norm{a}=\sqrt2$, $\norm{b}=\sqrt{18}$; $\cos(a,b)=6/\sqrt{36}=1$.
  \item $a\cdot c=1-1=0$, luego $\cos(a,c)=0$.
  \item Normalizando: $\hat a=\hat b=(1/\sqrt2,1/\sqrt2)$ y $\hat c=(1/\sqrt2,-1/\sqrt2)$. Entonces $\norm{\hat a-\hat b}=0$ y $\norm{\hat a-\hat c}=\norm{(0,\sqrt2)}=\sqrt2$.
\end{enumerate}
\resultado \textbf{$\norm{a-b}\approx 2{,}828$, $\norm{a-c}=2$, $\cos(a,b)=1$, $\cos(a,c)=0$.} La distancia elige $c$ y el coseno elige $b$: $b$ tiene la dirección de $a$, pero es tres veces más largo. Tras normalizar, ambas medidas eligen $b$.
\comprobacion la \cref{prop:coseno-distancia} da $\norm{\hat a-\hat c}^{2}=2-2\cdot 0=2$, que coincide con $(\sqrt2)^{2}$.
\end{solucion}

\begin{solucion}{eje:2-normalizar}
\datos $x=(3,4)$; $r=(1,0)$.
\planteamiento $\hat x=x/\norm{x}$; después evaluar ambos miembros de la identidad.
\pasos
\begin{enumerate}[label=\arabic*.]
  \item $\norm{x}=\sqrt{9+16}=5$ y $p=\hat x=(0{,}6;\,0{,}8)$.
  \item $p-r=(-0{,}4;\,0{,}8)$ y $\norm{p-r}^{2}=0{,}16+0{,}64=0{,}80$.
  \item Como $\norm{p}=\norm{r}=1$, $\cos(p,r)=p\cdot r=0{,}6$.
  \item $2-2\cos(p,r)=2-1{,}2=0{,}80$.
\end{enumerate}
\resultado \textbf{$\hat x=(0{,}6;\,0{,}8)$, y ambos miembros valen $0{,}80$.}
\comprobacion $\norm{p}^{2}=0{,}36+0{,}64=1$: $p$ es unitario y la identidad es aplicable.
\end{solucion}

\begin{solucion}{eje:2-tres}
\datos vectores de $\R^{2}$ con todos los productos escalares por pares negativos.
\planteamiento producto escalar negativo equivale a ángulo mayor que $90^\circ$ (\cref{tab:signo}); tres ángulos de $120^\circ$ completan la vuelta.
\pasos
\begin{enumerate}[label=\arabic*.]
  \item Se toman $u=(1,0)$, $v=(-1/2,\sqrt3/2)$ y $w=(-1/2,-\sqrt3/2)$, que son unitarios: $\norm{v}^{2}=\tfrac14+\tfrac34=1$, y lo mismo $w$.
  \item $u\cdot v=-\tfrac12$; $u\cdot w=-\tfrac12$; $v\cdot w=\tfrac14-\tfrac34=-\tfrac12$.
  \item Como son unitarios, $\cos\theta=-\tfrac12$ en los tres pares, luego $\theta=120^\circ$.
  \item (b) Supóngase que cuatro vectores no nulos tienen los seis productos negativos. Se ordenan por su ángulo polar y se llaman $g_1,\dots,g_4$ los ángulos que separan cada vector del siguiente en sentido antihorario, de modo que $g_1+g_2+g_3+g_4=360^\circ$.
  \item El ángulo entre dos vectores consecutivos es $\min(g_i,360^\circ-g_i)\le g_i$, y debe ser mayor que $90^\circ$; luego $g_i>90^\circ$ para todo $i$ y la suma supera $360^\circ$, lo cual es imposible.
\end{enumerate}
\resultado \textbf{(a) Sí: tres vectores unitarios a $120^\circ$, con los tres productos iguales a $-1/2$. (b) No: cuatro vectores del plano no pueden tener los seis productos negativos.}
\comprobacion $u+v+w=(1-\tfrac12-\tfrac12,\ 0+\tfrac{\sqrt3}{2}-\tfrac{\sqrt3}{2})=(0,0)$: las tres flechas se equilibran, como tres fuerzas iguales a $120^\circ$. En $\R^{3}$, en cambio, los cuatro vértices de un tetraedro regular centrado en el origen sí tienen los seis productos negativos.
\end{solucion}

\begin{solucion}{eje:2-distancia}
\datos $\norm{u}=\norm{v}=1$.
\planteamiento desarrollar la norma al cuadrado mediante el producto escalar.
\pasos
\begin{enumerate}[label=\arabic*.]
  \item $\norm{u-v}^{2}=(u-v)\cdot(u-v)=\norm{u}^{2}-2\,u\cdot v+\norm{v}^{2}$, por la linealidad y la simetría.
  \item Con normas unitarias: $\norm{u-v}^{2}=2-2\,u\cdot v$.
  \item También $\cos(u,v)=u\cdot v/(1\cdot 1)=u\cdot v$, luego $\norm{u-v}^{2}=2-2\cos(u,v)$.
  \item La función $c\mapsto 2-2c$ es estrictamente decreciente. Si $\cos(q,a)>\cos(q,b)$, entonces $\norm{q-a}^{2}<\norm{q-b}^{2}$ y, tomando raíces, $\norm{q-a}<\norm{q-b}$.
  \item Si $\cos(u,v)=-1$: $\norm{u-v}^{2}=4$ y $\norm{u-v}=2$.
\end{enumerate}
\resultado \textbf{$\norm{u-v}^{2}=2-2\cos(u,v)$; entre vectores unitarios, el más cercano por distancia es el de mayor coseno. Con $\cos=-1$, la distancia es~2 (puntos antipodales de la esfera unidad).}
\comprobacion el ejercicio~\ref{eje:2-normalizar} da $0{,}80=2-2\cdot 0{,}6$; con $\cos=1$ la distancia es~0, como debe ser para $u=v$.
\end{solucion}

% ---------------------------------------------------------------------
\section{Tema 3}

\begin{solucion}{eje:3-lineal}
\datos $f(x)=(x_1+x_2,3x_1)$, $g(x)=(x_1+1,x_2)$, $h(x)=(x_1x_2,x_2)$; $W\in\R^{5\times 3}$.
\planteamiento una aplicación lineal cumple $T(0)=0$ y $T(cx)=cT(x)$; si es lineal, su matriz tiene por columnas $T(e_1)$ y $T(e_2)$ (\cref{teo:representacion}).
\pasos
\begin{enumerate}[label=\arabic*.]
  \item $f$ es de la forma $Ax$ con $A=\begin{pmatrix}1&1\\3&0\end{pmatrix}$: $Ax=(x_1+x_2,\,3x_1)$. Por la \cref{prop:matriz-lineal}, es lineal. Columnas: $f(e_1)=(1,3)$ y $f(e_2)=(1,0)$.
  \item $g(0)=(1,0)\neq 0$, luego $g$ no es lineal.
  \item $h(2,2)=(4,2)$ y $2h(1,1)=2(1,1)=(2,2)$; como difieren, $h$ no es lineal.
  \item (b) $W$ actúa sobre vectores de $\R^{3}$ (tres columnas) y produce vectores de $\R^{5}$ (cinco filas).
\end{enumerate}
\resultado \textbf{Solo $f$ es lineal, con matriz $\begin{pmatrix}1&1\\3&0\end{pmatrix}$. $W$ transforma $\R^{3}$ en $\R^{5}$.}
\comprobacion $f(1,1)=(2,3)$ y $A(1,1)=(1+1,\,3+0)=(2,3)$.
\end{solucion}

\begin{solucion}{eje:3-dos-lecturas}
\datos $M=\begin{pmatrix}1&0&2\\2&1&0\\0&3&1\end{pmatrix}$, $x=(2,-1,3)$.
\planteamiento las dos lecturas de la \cref{prop:filas-columnas}.
\pasos
\begin{enumerate}[label=\arabic*.]
  \item Fila 1: $1\cdot 2+0\cdot(-1)+2\cdot 3=8$.
  \item Fila 2: $2\cdot 2+1\cdot(-1)+0\cdot 3=3$.
  \item Fila 3: $0\cdot 2+3\cdot(-1)+1\cdot 3=0$.
  \item Columnas: $2(1,2,0)+(-1)(0,1,3)+3(2,0,1)=(2,4,0)+(0,-1,-3)+(6,0,3)=(8,3,0)$.
\end{enumerate}
\resultado \textbf{$Mx=(8,3,0)$.} La tercera fila de $M$ es ortogonal a $x$, porque su producto escalar es~0.
\comprobacion las dos lecturas dan el mismo vector, como garantiza la \cref{prop:filas-columnas}.
\end{solucion}

\begin{solucion}{eje:3-base}
\datos $T(e_1)=(1,2)$, $T(e_2)=(-1,1)$; $x=(3,-2)$.
\planteamiento por el \cref{teo:representacion}, las columnas de la matriz son $T(e_1)$ y $T(e_2)$.
\pasos
\begin{enumerate}[label=\arabic*.]
  \item $A=\begin{pmatrix}1&-1\\2&1\end{pmatrix}$.
  \item Con la matriz: $Ax=(1\cdot 3+(-1)(-2),\ 2\cdot 3+1\cdot(-2))=(5,4)$.
  \item Con la linealidad: $T(3e_1-2e_2)=3T(e_1)-2T(e_2)=(3,6)+(2,-2)=(5,4)$.
\end{enumerate}
\resultado \textbf{$A=\begin{pmatrix}1&-1\\2&1\end{pmatrix}$ y $T(3,-2)=(5,4)$.}
\comprobacion $Ae_1=(1,2)$ y $Ae_2=(-1,1)$, que son los datos.
\end{solucion}

\begin{solucion}{eje:3-rotacion}
\datos $\theta=60^\circ$, $\cos 60^\circ=\tfrac12$, $\sin 60^\circ=\tfrac{\sqrt3}{2}$.
\planteamiento sustituir en la matriz $R_\theta$ del \cref{tab:transformaciones} y multiplicar.
\pasos
\begin{enumerate}[label=\arabic*.]
  \item $R_{60^\circ}=\begin{pmatrix}1/2&-\sqrt3/2\\ \sqrt3/2&1/2\end{pmatrix}$.
  \item $R_{60^\circ}(2,0)=(1,\sqrt3)$, cuya norma es $\sqrt{1+3}=2=\norm{(2,0)}$.
  \item $R_{60^\circ}(1,1)=\bigl(\tfrac{1-\sqrt3}{2},\tfrac{\sqrt3+1}{2}\bigr)\approx(-0{,}3660;\,1{,}366)$.
  \item Norma al cuadrado: $\tfrac14\bigl[(1-\sqrt3)^{2}+(1+\sqrt3)^{2}\bigr]=\tfrac14\bigl[(4-2\sqrt3)+(4+2\sqrt3)\bigr]=2$, luego la norma es $\sqrt2=\norm{(1,1)}$.
\end{enumerate}
\resultado \textbf{$R_{60^\circ}(2,0)=(1,\sqrt3)$ y $R_{60^\circ}(1,1)\approx(-0{,}3660;\,1{,}366)$; las normas ($2$ y $\sqrt2$) se conservan.}
\comprobacion el vector $(1,\sqrt3)$ forma con el eje $x$ un ángulo de $\arctan\sqrt3=60^\circ$, como se esperaba.
\end{solucion}

\begin{solucion}{eje:3-markov}
\datos $A=\begin{pmatrix}0{,}8&0{,}3\\0{,}2&0{,}7\end{pmatrix}$, $u_0=(1,0)$, $s=(0{,}6;\,0{,}4)$.
\planteamiento aplicar $A$ repetidamente; para (c), intercambiar el orden de las sumas.
\pasos
\begin{enumerate}[label=\arabic*.]
  \item $u_1=Au_0=(0{,}8;\,0{,}2)$, la primera columna de $A$.
  \item $u_2=Au_1=(0{,}8\cdot 0{,}8+0{,}3\cdot 0{,}2;\ 0{,}2\cdot 0{,}8+0{,}7\cdot 0{,}2)=(0{,}64+0{,}06;\ 0{,}16+0{,}14)=(0{,}70;\,0{,}30)$.
  \item $u_3=Au_2=(0{,}8\cdot 0{,}7+0{,}3\cdot 0{,}3;\ 0{,}2\cdot 0{,}7+0{,}7\cdot 0{,}3)=(0{,}56+0{,}09;\ 0{,}14+0{,}21)=(0{,}65;\,0{,}35)$.
  \item Las componentes de $u_0,\dots,u_3$ son no negativas y suman~1.
  \item (b) $As=(0{,}48+0{,}12;\ 0{,}12+0{,}28)=(0{,}6;\,0{,}4)=s$.
  \item (c) $\unos^{\top}(Au)=\sum_{i}\sum_{j}A_{ij}u_j=\sum_ju_j\Bigl(\sum_iA_{ij}\Bigr)=\sum_ju_j\cdot 1=\unos^{\top}u$.
\end{enumerate}
\resultado \textbf{$u_1=(0{,}8;\,0{,}2)$, $u_2=(0{,}7;\,0{,}3)$, $u_3=(0{,}65;\,0{,}35)$, todos con suma~1; $As=s$; una matriz con columnas de suma~1 conserva la suma de las componentes.}
\comprobacion los $u_k$ se acercan a $s$: la primera componente pasa de $1$ a $0{,}8$, $0{,}7$ y $0{,}65$, y la distancia a $0{,}6$ se reduce a la mitad en cada paso ($0{,}4$; $0{,}2$; $0{,}1$; $0{,}05$).
\end{solucion}

\begin{solucion}{eje:3-rectas}
\datos $T\colon\R^{n}\to\R^{m}$ lineal; $p,d\in\R^{n}$; $b\neq 0$.
\planteamiento aplicar la \cref{def:lineal} a combinaciones concretas.
\pasos
\begin{enumerate}[label=\arabic*.]
  \item (a) Por linealidad, $T(p+td)=T(p)+t\,T(d)$ para todo $t$. Si $T(d)\neq 0$, la imagen es la recta que pasa por $T(p)$ con dirección $T(d)$; si $T(d)=0$, la imagen es el punto $T(p)$.
  \item (b) Con $a=b=\tfrac12$ en la \cref{def:lineal}: $T\bigl(\tfrac12u+\tfrac12v\bigr)=\tfrac12T(u)+\tfrac12T(v)$.
  \item (c) Sea $S(x)=x+b$. La imagen de la recta es $\{(p+b)+td\}$, otra recta, y $S\bigl(\tfrac12(u+v)\bigr)=\tfrac12(u+v)+b=\tfrac12\bigl((u+b)+(v+b)\bigr)$.
  \item Sin embargo, $S(0)=b\neq 0$, mientras que toda aplicación lineal fija el origen. Luego $S$ no es lineal.
\end{enumerate}
\resultado \textbf{Una aplicación lineal transforma rectas en rectas (o en un punto, si $T(d)=0$) y conserva puntos medios; las traslaciones tienen ambas propiedades sin ser lineales.} Las aplicaciones de la forma $x\mapsto Ax+b$ se llaman \emph{afines}~\cite[§8.1]{vmls}; una capa de red neuronal con sesgo es afín (módulo~9).
\comprobacion la proyección sobre el eje $x$ del \cref{tab:transformaciones} envía la recta vertical $\{(1,0)+t(0,1)\}$ al punto $(1,0)$, porque $T(0,1)=(0,0)$.
\end{solucion}

% ---------------------------------------------------------------------
\section{Tema 4}

\begin{solucion}{eje:4-formas}
\datos tamaños de las matrices y $u,v\in\R^{3}$.
\planteamiento $(m\times n)(n\times p)=m\times p$, solo si coinciden las dimensiones interiores.
\pasos
\begin{enumerate}[label=(\alph*)]
  \item $(3\times 4)(4\times 2)=3\times 2$. El producto $(2\times 3)(2\times 3)$ no está definido, porque $3\neq 2$.
  \item $u^{\top}v$ es $(1\times 3)(3\times 1)=1\times 1$, un número (el producto escalar); $uv^{\top}$ es $(3\times 1)(1\times 3)=3\times 3$, un producto exterior.
  \item \textbf{Falso.} En el \cref{ej:intercambio}, $AB\neq BA$.
\end{enumerate}
\resultado \textbf{(a) $3\times 2$; no. (b) $1\times 1$ y $3\times 3$. (c) Falso.}
\comprobacion con $u=v=(1,1,1)$: $u^{\top}v=3$, y $uv^{\top}$ es la matriz $3\times 3$ de unos, cuya traza es también~3.
\end{solucion}

\begin{solucion}{eje:4-rs}
\datos $R=\begin{pmatrix}0&-1\\1&0\end{pmatrix}$, $S=\begin{pmatrix}2&0\\0&0{,}5\end{pmatrix}$, $x=(3,1)$.
\planteamiento entrada $(i,j)$ del producto $=$ fila $i$ del primer factor por columna $j$ del segundo.
\pasos
\begin{enumerate}[label=\arabic*.]
  \item $RS$: fila 1, $(0\cdot 2+(-1)\cdot 0;\ 0\cdot 0+(-1)\cdot 0{,}5)=(0;\,-0{,}5)$; fila 2, $(1\cdot 2+0;\ 0+0)=(2;\,0)$.
  \item $SR$: fila 1, $(2\cdot 0+0\cdot 1;\ 2\cdot(-1)+0\cdot 0)=(0;\,-2)$; fila 2, $(0\cdot 0+0{,}5\cdot 1;\ 0+0)=(0{,}5;\,0)$.
  \[
  RS=\begin{pmatrix}0&-0{,}5\\2&0\end{pmatrix},\qquad SR=\begin{pmatrix}0&-2\\0{,}5&0\end{pmatrix}.
  \]
  \item $RSx=(0\cdot 3-0{,}5\cdot 1;\ 2\cdot 3+0)=(-0{,}5;\,6)$. Por pasos: $Sx=(6;\,0{,}5)$ y $R(6;\,0{,}5)=(-0{,}5;\,6)$.
  \item $SRx=(0\cdot 3-2\cdot 1;\ 0{,}5\cdot 3+0)=(-2;\,1{,}5)$. Por pasos: $Rx=(-1,3)$ y $S(-1,3)=(-2;\,1{,}5)$.
\end{enumerate}
\resultado \textbf{$RS=\begin{pmatrix}0&-0{,}5\\2&0\end{pmatrix}$, $SR=\begin{pmatrix}0&-2\\0{,}5&0\end{pmatrix}$; $RSx=(-0{,}5;\,6)$ y $SRx=(-2;\,1{,}5)$.} Los resultados coinciden con la aplicación sucesiva, como afirma el \cref{teo:composicion}.
\comprobacion $R$ conserva normas: $\norm{Sx}^{2}=36+0{,}25=36{,}25$ y $\norm{R(Sx)}^{2}=0{,}25+36=36{,}25$.
\end{solucion}

\begin{solucion}{eje:4-ab-ba}
\datos $A\in\R^{2\times 3}$ y $B\in\R^{3\times 2}$ del enunciado.
\planteamiento $(2\times 3)(3\times 2)=2\times 2$ y $(3\times 2)(2\times 3)=3\times 3$.
\pasos
\begin{enumerate}[label=\arabic*.]
  \item $(AB)_{11}=1\cdot 1+0\cdot 0+2\cdot 3=7$; $(AB)_{12}=1\cdot 2+0\cdot 1+2\cdot 0=2$.
  \item $(AB)_{21}=0\cdot 1+1\cdot 0+(-1)\cdot 3=-3$; $(AB)_{22}=0\cdot 2+1\cdot 1+(-1)\cdot 0=1$.
  \item Filas de $BA$ (\cref{prop:cuatro}(ii)): $(1,2)A=(1,2,0)$; $(0,1)A=(0,1,-1)$; $(3,0)A=(3,0,6)$.
  \[
  AB=\begin{pmatrix}7&2\\-3&1\end{pmatrix},\qquad BA=\begin{pmatrix}1&2&0\\0&1&-1\\3&0&6\end{pmatrix}.
  \]
  \item $AA$ sería $(2\times 3)(2\times 3)$, que no está definido.
\end{enumerate}
\resultado \textbf{$AB$ es $2\times 2$ y $BA$ es $3\times 3$, con los valores indicados; $AA$ no existe.}
\comprobacion las trazas de $AB$ y de $BA$ coinciden siempre: $7+1=8$ y $1+1+6=8$.
\end{solucion}

\begin{solucion}{eje:4-distributiva}
\datos $A=\begin{pmatrix}2&1\\1&3\end{pmatrix}$, $B=\begin{pmatrix}1&0\\0&0\end{pmatrix}$, $C=\begin{pmatrix}0&0\\2&1\end{pmatrix}$.
\planteamiento calcular cada producto entrada a entrada.
\pasos
\begin{enumerate}[label=\arabic*.]
  \item $AB=\begin{pmatrix}2&0\\1&0\end{pmatrix}$ y $AC=\begin{pmatrix}2\cdot 0+1\cdot 2&2\cdot 0+1\cdot 1\\1\cdot 0+3\cdot 2&1\cdot 0+3\cdot 1\end{pmatrix}=\begin{pmatrix}2&1\\6&3\end{pmatrix}$.
  \item $AB+AC=\begin{pmatrix}4&1\\7&3\end{pmatrix}$.
  \item $B+C=\begin{pmatrix}1&0\\2&1\end{pmatrix}$ y $A(B+C)=\begin{pmatrix}2+2&0+1\\1+6&0+3\end{pmatrix}=\begin{pmatrix}4&1\\7&3\end{pmatrix}$.
  \item $BA=\begin{pmatrix}2&1\\0&0\end{pmatrix}$ y $CA=\begin{pmatrix}0&0\\2\cdot 2+1\cdot 1&2\cdot 1+1\cdot 3\end{pmatrix}=\begin{pmatrix}0&0\\5&5\end{pmatrix}$.
  \item $BA+CA=\begin{pmatrix}2&1\\5&5\end{pmatrix}$.
\end{enumerate}
\resultado \textbf{$AB+AC=A(B+C)=\begin{pmatrix}4&1\\7&3\end{pmatrix}$, mientras que $BA+CA=(B+C)A=\begin{pmatrix}2&1\\5&5\end{pmatrix}$ es distinta.} Cada distributiva vale por su lado, pero no se puede cambiar el lado.
\comprobacion $(B+C)A=\begin{pmatrix}1&0\\2&1\end{pmatrix}\begin{pmatrix}2&1\\1&3\end{pmatrix}=\begin{pmatrix}2&1\\4+1&2+3\end{pmatrix}=\begin{pmatrix}2&1\\5&5\end{pmatrix}$, que coincide con el paso~5.
\end{solucion}

\begin{solucion}{eje:4-coste}
\datos $A\in\R^{3\times 5}$, $B\in\R^{5\times 2}$, $C\in\R^{2\times 4}$; coste $mnp$ (\cref{prop:coste}).
\planteamiento sumar el coste de los dos productos de cada orden.
\pasos
\begin{enumerate}[label=\arabic*.]
  \item $(AB)C$: $AB$ cuesta $3\cdot 5\cdot 2=30$ y es $3\times 2$; $(AB)C$ cuesta $3\cdot 2\cdot 4=24$. Total: $54$.
  \item $A(BC)$: $BC$ cuesta $5\cdot 2\cdot 4=40$ y es $5\times 4$; $A(BC)$ cuesta $3\cdot 5\cdot 4=60$. Total: $100$.
  \item (b) $(u^{\top}v)w^{\top}$: el escalar $u^{\top}v$ cuesta $N$ y multiplicarlo por $w^{\top}$ cuesta $N$. Total: $2N$.
  \item $u^{\top}(vw^{\top})$: $vw^{\top}$ es $N\times N$ y cuesta $N^{2}$; $u^{\top}(vw^{\top})$ cuesta otras $N^{2}$. Total: $2N^{2}$.
  \item Con $N=4096$: $2N=8192$ y $2N^{2}=33\,554\,432$, un factor $N=4096$ mayor.
\end{enumerate}
\resultado \textbf{$(AB)C$ cuesta 54 multiplicaciones y $A(BC)$ cuesta 100. $(u^{\top}v)w^{\top}$ cuesta $2N$ y $u^{\top}(vw^{\top})$ cuesta $2N^{2}$.} El resultado es el mismo por la asociativa, pero el coste no.
\comprobacion ambos órdenes de (a) producen una matriz $3\times 4$; en (b), ambos producen el vector fila $(u^{\top}v)\,w^{\top}$ de $N$ componentes.
\end{solucion}

\begin{solucion}{eje:4-proyeccion}
\datos $A=\begin{pmatrix}1&0\\0&1\\0&0\end{pmatrix}$, $b=(4,-1,2)$, $a=(1,1,1)$.
\planteamiento calcular $P=AA^{\top}$ (filas de $A$ por filas de $A$) y aplicar $P^{2}=P$; en (c), $P_a=aa^{\top}/(a^{\top}a)$.
\pasos
\begin{enumerate}[label=\arabic*.]
  \item (a) $(AA^{\top})_{ij}$ es el producto escalar de las filas $i$ y $j$ de $A$, que son $(1,0)$, $(0,1)$ y $(0,0)$: $P=\diag(1,1,0)$.
  \item $Pb=(4,-1,0)$: se conservan las dos primeras coordenadas y se anula la tercera.
  \item (b) $P^{2}=\diag(1^{2},1^{2},0^{2})=P$.
  \item $(I-P)^{2}=I-2P+P^{2}=I-2P+P=I-P$, usando $P^{2}=P$.
  \item $I-P=\diag(0,0,1)$ y $(I-P)b=(0,0,2)$: proyecta sobre el eje $z$, perpendicular al plano generado por las columnas de $A$.
  \item (c) $a^{\top}a=3$ y $aa^{\top}$ es la matriz $3\times 3$ de unos; $P_a=\tfrac13 aa^{\top}$.
  \item $P_ab=\tfrac13(4-1+2)(1,1,1)=\tfrac53(1,1,1)\approx(1{,}667;\,1{,}667;\,1{,}667)$.
\end{enumerate}
\resultado \textbf{$P=\diag(1,1,0)$ y $Pb=(4,-1,0)$; $P^{2}=P$ y $(I-P)^{2}=I-P$; $I-P$ proyecta sobre el eje $z$; $P_ab=\tfrac53(1,1,1)$, como en el \cref{ej:proyeccion}.}
\comprobacion $Pb+(I-P)b=(4,-1,0)+(0,0,2)=b$ y $(4,-1,0)\cdot(0,0,2)=0$: $b$ se descompone en dos piezas ortogonales.
\end{solucion}

\begin{solucion}{eje:4-ortogonal}
\datos $R_\theta=\begin{pmatrix}\cos\theta&-\sin\theta\\ \sin\theta&\cos\theta\end{pmatrix}$; fórmulas del coseno y del seno de la suma.
\planteamiento calcular los productos entrada a entrada; para (c), usar $(AB)^{\top}=B^{\top}A^{\top}$.
\pasos
\begin{enumerate}[label=\arabic*.]
  \item (a) Las columnas de $R_\theta$ son $c_1=(\cos\theta,\sin\theta)$ y $c_2=(-\sin\theta,\cos\theta)$. $(R_\theta^{\top}R_\theta)_{ij}=c_i\cdot c_j$: $c_1\cdot c_1=c_2\cdot c_2=\cos^{2}\theta+\sin^{2}\theta=1$ y $c_1\cdot c_2=-\cos\theta\sin\theta+\sin\theta\cos\theta=0$. Luego $R_\theta^{\top}R_\theta=I_2$.
  \item (b) La primera columna de $R_\theta R_\varphi$ es
  \[
  R_\theta\begin{pmatrix}\cos\varphi\\ \sin\varphi\end{pmatrix}=\begin{pmatrix}\cos\theta\cos\varphi-\sin\theta\sin\varphi\\ \sin\theta\cos\varphi+\cos\theta\sin\varphi\end{pmatrix}=\begin{pmatrix}\cos(\theta+\varphi)\\ \sin(\theta+\varphi)\end{pmatrix}.
  \]
  \item La segunda columna es
  \[
  R_\theta\begin{pmatrix}-\sin\varphi\\ \cos\varphi\end{pmatrix}=\begin{pmatrix}-\cos\theta\sin\varphi-\sin\theta\cos\varphi\\ -\sin\theta\sin\varphi+\cos\theta\cos\varphi\end{pmatrix}=\begin{pmatrix}-\sin(\theta+\varphi)\\ \cos(\theta+\varphi)\end{pmatrix}.
  \]
  Por tanto, $R_\theta R_\varphi=R_{\theta+\varphi}$.
  \item Como $\theta+\varphi=\varphi+\theta$, $R_\theta R_\varphi=R_{\theta+\varphi}=R_\varphi R_\theta$.
  \item (c) Si $Q_1^{\top}Q_1=Q_2^{\top}Q_2=I$: $(Q_1Q_2)^{\top}(Q_1Q_2)=Q_2^{\top}Q_1^{\top}Q_1Q_2=Q_2^{\top}Q_2=I$.
  \item (d) Para $D=\diag(\lambda_1,\lambda_2)$ y $R=R_{90^\circ}$: $RD=\begin{pmatrix}0&-\lambda_2\\ \lambda_1&0\end{pmatrix}$ y $DR=\begin{pmatrix}0&-\lambda_1\\ \lambda_2&0\end{pmatrix}$, que coinciden si y solo si $\lambda_1=\lambda_2$.
\end{enumerate}
\resultado \textbf{Las rotaciones son ortogonales, se componen sumando ángulos y conmutan entre sí; el producto de ortogonales es ortogonal.} $S$ no conmuta con $R$ porque escala los dos ejes por factores distintos ($2\neq 0{,}5$) y no es ortogonal.
\comprobacion con $\theta=\varphi=90^\circ$: $R_{90^\circ}^{2}=\begin{pmatrix}0&-1\\1&0\end{pmatrix}^{2}=\begin{pmatrix}-1&0\\0&-1\end{pmatrix}=R_{180^\circ}$; con los datos del paso~6, $\lambda_1=2$ y $\lambda_2=0{,}5$ reproducen $RS$ y $SR$ del ejercicio~\ref{eje:4-rs}.
\end{solucion}

% ---------------------------------------------------------------------
\section{Tema 5}

\begin{solucion}{eje:5-tamano}
\datos $K=50\,257$; $d=768$ o $d=1600$; 4 bytes por número.
\planteamiento $E$ tiene una fila por token y una columna por dimensión: $Kd$ parámetros.
\pasos
\begin{enumerate}[label=\arabic*.]
  \item Tamaño: $E\in\R^{50\,257\times 768}$.
  \item Parámetros: $50\,257\cdot 768=50\,257\cdot 700+50\,257\cdot 68=35\,179\,900+3\,417\,476=38\,597\,376$.
  \item Memoria: $38\,597\,376\cdot 4=154\,389\,504$ bytes $\approx 154$~MB.
  \item Con $d=1600$: $50\,257\cdot 1600=80\,411\,200$ parámetros y $80\,411\,200\cdot 4=321\,644\,800$ bytes $\approx 322$~MB.
\end{enumerate}
\resultado \textbf{$E$ es $50\,257\times 768$, con $38\,597\,376$ parámetros ($\approx 154$~MB); con $d=1600$, $80\,411\,200$ parámetros ($\approx 322$~MB).}
\comprobacion orden de magnitud: $5\cdot 10^{4}\cdot 8\cdot 10^{2}=4\cdot 10^{7}\approx 3{,}9\cdot 10^{7}$.
\end{solucion}

\begin{solucion}{eje:5-onehot}
\datos $e_i\in\R^{K}$ one-hot; $E\in\R^{K\times d}$; en GPT-2, $K=50\,257$ y $d=768$.
\planteamiento escribir la componente $j$ del producto con la \cref{def:producto}.
\pasos
\begin{enumerate}[label=\arabic*.]
  \item $(e_i^{\top}E)_j=\sum_{k}(e_i)_kE_{kj}$.
  \item $(e_i)_k=0$ si $k\neq i$ y $(e_i)_i=1$, luego solo sobrevive el término $k=i$: $(e_i^{\top}E)_j=E_{ij}$.
  \item Esto vale para $j=1,\dots,d$: $e_i^{\top}E=(E_{i1},\dots,E_{id})$, la fila $i$.
  \item El producto requiere $Kd$ multiplicaciones; la lectura copia $d$ números. El cociente es $Kd/d=K=50\,257$.
\end{enumerate}
\resultado \textbf{$e_i^{\top}E$ es la fila $i$ de $E$; el producto requiere unas $50\,000$ veces más operaciones que la lectura directa.}
\comprobacion en el \cref{ej:embedding}, $(0,0,1,0)\,E=(0{,}5;\,0{,}5;\,0{,}5)$, que es la fila~2.
\end{solucion}

\begin{solucion}{eje:5-busqueda}
\datos $\hat E$ con filas unitarias; $K=50\,257$, $d=768$; lote de $B$ consultas.
\planteamiento $\cos(v_w,q)=\hat v_w\cdot\hat q$ y la lectura por filas del producto matriz-vector (\cref{prop:filas-columnas}).
\pasos
\begin{enumerate}[label=\arabic*.]
  \item (a) La componente $w$ de $\hat E\hat q$ es $\hat E_{w,:}\cdot\hat q=\cos(v_w,q)$. Un único producto matriz-vector da los $K$ cosenos. Como dividir por $\norm{q}>0$ no altera el orden (\cref{prop:coseno}(i)), basta calcular $\hat Eq$.
  \item (b) $\hat Eq$ requiere $Kd=50\,257\cdot 768=38\,597\,376$ multiplicaciones. La normalización de $E$ se hace una sola vez y no se repite en cada consulta.
  \item (c) Se apilan las consultas como columnas de $Q\in\R^{d\times B}$ y se calcula $\hat EQ\in\R^{K\times B}$, con $KdB$ multiplicaciones en un solo producto de matrices, que la GPU ejecuta eficientemente.
  \item Para vocabularios o colecciones mucho mayores se usan índices de búsqueda aproximada de vecinos, que evitan comparar con todos los candidatos (módulo~22).
\end{enumerate}
\resultado \textbf{Con filas normalizadas de antemano, los $K$ cosenos son $\hat Eq$ salvo un factor positivo común; cuesta unos $3{,}86\cdot 10^{7}$ multiplicaciones por consulta, y un lote se procesa con el producto $\hat EQ$.}
\comprobacion el coste coincide con el del producto $e_i^{\top}E$ del ejercicio~\ref{eje:5-onehot}, porque ambos recorren una vez todas las entradas de una matriz $K\times d$.
\end{solucion}

\begin{solucion}{eje:5-analogias}
\datos vectores del \cref{ej:juguete} y los cinco vectores geográficos del enunciado.
\planteamiento aritmética componente a componente y similitud coseno.
\pasos
\begin{enumerate}[label=\arabic*.]
  \item (a) $v_{\text{hombre}}\cdot v_{\text{mujer}}=0+(-1)+4+0=3$ y ambas normas valen $\sqrt5$: $\cos=3/5=0{,}6$.
  \item (b) $v_{\text{rey}}-v_{\text{hombre}}=(2,0,0,0)$; sumando $v_{\text{mujer}}$: $(2,-1,2,0)=v_{\text{reina}}$.
  \item (c) $v_{\text{Madrid}}-v_{\text{España}}=(-1,1,0,0)$; sumando $v_{\text{Francia}}$: $q=(0,1,0,1)=v_{\text{París}}$, luego $\cos(q,v_{\text{París}})=1$.
  \item $q\cdot v_{\text{Lyon}}=0{,}5+1=1{,}5$; $\norm{q}=\sqrt2\approx 1{,}414$; $\norm{v_{\text{Lyon}}}=\sqrt{1{,}25}\approx 1{,}118$.
  \item $\cos(q,v_{\text{Lyon}})=1{,}5/\sqrt{2{,}5}=1{,}5/1{,}581\approx 0{,}9487$.
\end{enumerate}
\resultado \textbf{(a) $0{,}6$; (b) se obtiene exactamente «reina»; (c) $q=(0,1,0,1)$, con coseno~1 con París y $\approx 0{,}9487$ con Lyon.} París gana, pero Lyon queda cerca.
\comprobacion $0{,}9487<1$, como exige el \cref{teo:cs}; y (b) es el caso exacto de la \cref{prop:paralelogramo}.
\end{solucion}

\begin{solucion}{eje:5-secuencia}
\datos $E\in\R^{4\times 3}$ del \cref{ej:embedding}; secuencia $(2,0,2)$.
\planteamiento \cref{prop:onehot}: cada fila de $O$ es el one-hot de un token y $X=OE$.
\pasos
\begin{enumerate}[label=\arabic*.]
  \item (a) Columna 1: $0\cdot 0{,}1+0\cdot 1+1\cdot 0{,}5+0\cdot(-0{,}2)=0{,}5$. Las columnas 2 y 3 dan también $0{,}5$.
  \item (b) $O=\begin{pmatrix}0&0&1&0\\1&0&0&0\\0&0&1&0\end{pmatrix}\in\R^{3\times 4}$.
  \item $X=OE$: sus filas son las filas 2, 0 y 2 de $E$:
  \[
  X=\begin{pmatrix}0{,}5&0{,}5&0{,}5\\0{,}1&0{,}2&0\\0{,}5&0{,}5&0{,}5\end{pmatrix}\in\R^{3\times 3}.
  \]
\end{enumerate}
\resultado \textbf{(a) $e_2^{\top}E=(0{,}5;\,0{,}5;\,0{,}5)$. (b) $X$ es $3\times 3$ ($n\times d$), con las filas primera y tercera iguales.}
\comprobacion tamaños: $(3\times 4)(4\times 3)=3\times 3$.
\end{solucion}

\begin{solucion}{eje:5-logits}
\datos $E$ del \cref{ej:embedding}; $h=(1,0,-1)$.
\planteamiento $z_i=E_{i,:}\cdot h$ y $\softmax(z)_i=e^{z_i}/\sum_je^{z_j}$.
\pasos
\begin{enumerate}[label=\arabic*.]
  \item $z_0=0{,}1-0=0{,}1$; $z_1=1+1=2$; $z_2=0{,}5-0{,}5=0$; $z_3=-0{,}2-0{,}1=-0{,}3$.
  \item Exponenciales: $e^{0{,}1}=1{,}105$; $e^{2}=7{,}389$; $e^{0}=1$; $e^{-0{,}3}=0{,}7408$. Suma: $10{,}24$.
  \item Softmax: $1{,}105/10{,}24=0{,}1080$; $7{,}389/10{,}24=0{,}7219$; $1/10{,}24=0{,}0977$; $0{,}7408/10{,}24=0{,}0724$.
  \item La fila~1, $(1,0,-1)$, es igual a $h$: $\cos(E_{1,:},h)=1$ y $z_1=\norm{E_{1,:}}\norm{h}\cos=\sqrt2\cdot\sqrt2\cdot 1=2$.
\end{enumerate}
\resultado \textbf{$z=(0{,}1;\,2;\,0;\,-0{,}3)$ y $\softmax(z)\approx(0{,}1080;\,0{,}7219;\,0{,}0977;\,0{,}0724)$. El token más probable es el~1, cuyo embedding apunta exactamente en la dirección de $h$.}
\comprobacion $0{,}1080+0{,}7219+0{,}0977+0{,}0724=1{,}0000$.
\end{solucion}

\begin{solucion}{eje:5-atencion}
\datos (a) $q=(1,0)$, $k_1=(1,1)$, $k_2=(0,2)$, $k_3=(-1,0)$, $d_k=2$. (b) Componentes independientes de media~0 y varianza~1. (c) $d=64$.
\planteamiento (a) puntuaciones y softmax; (b) linealidad de la esperanza y aditividad de la varianza para sumandos independientes.
\pasos
\begin{enumerate}[label=\arabic*.]
  \item (a) $q\cdot k_1=1$, $q\cdot k_2=0$, $q\cdot k_3=-1$; divididos por $\sqrt2\approx 1{,}414$: $(0{,}7071;\,0;\,-0{,}7071)$.
  \item $e^{0{,}7071}=2{,}028$, $e^{0}=1$, $e^{-0{,}7071}=0{,}4931$; suma $3{,}521$; softmax $(0{,}5760;\,0{,}2840;\,0{,}1400)$.
  \item (b) $q\cdot k=\sum_{i=1}^{d}q_ik_i$. Por independencia, $\E[q_ik_i]=\E[q_i]\E[k_i]=0$, luego $\E[q\cdot k]=0$.
  \item $\Var(q_ik_i)=\E[q_i^{2}k_i^{2}]-0^{2}=\E[q_i^{2}]\,\E[k_i^{2}]=1\cdot 1=1$, porque $\E[X^{2}]=\Var(X)+\E[X]^{2}=1$.
  \item Los $d$ sumandos $q_ik_i$ dependen de pares disjuntos de variables independientes, luego son independientes y sus varianzas se suman: $\Var(q\cdot k)=d$.
  \item (c) Con $d=64$, la desviación típica es $\sqrt{64}=8$. La softmax de $(8,0,-8)$ es $(2981;\,1;\,0{,}000335)/2982\approx(0{,}99966;\,0{,}00034;\,1{,}1\cdot 10^{-7})$: casi un vector one-hot.
  \item La softmax de $(1,0,-1)$ es $(2{,}718;\,1;\,0{,}3679)/4{,}086\approx(0{,}6652;\,0{,}2447;\,0{,}0900)$: una distribución repartida.
\end{enumerate}
\resultado \textbf{(a) puntuaciones $(0{,}7071;\,0;\,-0{,}7071)$ y pesos $(0{,}576;\,0{,}284;\,0{,}140)$. (b) $\E[q\cdot k]=0$ y $\Var(q\cdot k)=d$. (c) Sin dividir por $\sqrt d$ la softmax se satura; dividiendo, la distribución queda repartida y los gradientes no se anulan.}
\comprobacion (a) $0{,}576+0{,}284+0{,}140=1{,}000$. (b) Una simulación con $10^{5}$ pares de vectores gaussianos de dimensión 64 da una varianza muestral de $63{,}8$ antes de dividir por~8 y de $1{,}00$ después.
\end{solucion}
